% ============================================================
% pretty-charts TikZ 流程图预置样式
% 论文版与 PPT 版共用同一套节点形状，仅配色/线宽不同：
%   \pcPaperMode  ← 论文（默认）：黑白灰，无填充色
%   \pcPPTMode    ← 演示：showcase 色板彩色，粗线大字
% 用法：先 % ============================================================
% pretty-charts TikZ/pgfplots 导言区模板 —— academic 档（T1 出版级）
% 编译方式：XeLaTeX。business/showcase 档改字号与 cycle list（见文件末尾注释）
% ============================================================
\usepackage{pgfplots}
\pgfplotsset{compat=1.18}
\usetikzlibrary{arrows.meta, positioning, calc}

% ---- 图表字体：中文黑体 + 西文 Arial（XeLaTeX；正文仍可用全局宋体+TNR，互不影响）----
% 注意：Windows 的 Noto Sans SC 多为可变字体(VF)，xdvipdfmx 无法嵌入（fatal: Invalid font），
% 故 TikZ 导出 PDF 用 Microsoft YaHei；装有静态版思源黑体(Source Han Sans SC)的机器可换回。
\usepackage{fontspec}
\usepackage{xeCJK}
\newCJKfontfamily\chartcjk{Microsoft YaHei}
\newcommand{\chartfont}[1]{{\chartcjk\fontspec{Arial}#1}}

% ---- 色板：academic（Okabe-Ito 重排，与 assets/palettes/academic.json 一致）----
\definecolor{pcBlue}{HTML}{0072B2}
\definecolor{pcOrange}{HTML}{D55E00}
\definecolor{pcGreen}{HTML}{009E73}
\definecolor{pcMagenta}{HTML}{CC79A7}
\definecolor{pcYellow}{HTML}{E69F00}
\definecolor{pcCyan}{HTML}{56B4E9}
\definecolor{pcGray}{HTML}{999999}
\definecolor{pcGrid}{HTML}{CCCCCC}
\definecolor{pcAxis}{HTML}{333333}
\definecolor{pcTick}{HTML}{4D4D4D}

% ---- 全局轴风格：去顶右刺、浅色 y 网格、字号层级 ----
\pgfplotsset{
  every axis/.append style={
    axis line style={pcAxis, line width=0.4pt},
    label style={font=\small, text=pcAxis},
    tick label style={font=\footnotesize, text=pcTick},
    title style={font=\small\bfseries, text=pcTitleBlack},
    legend style={font=\footnotesize, draw=none, fill=none, nodes={scale=0.9, transform shape}},
    grid=major,
    major grid style={pcGrid, line width=0.4pt},
    axis x line*=bottom,
    axis y line*=left,
    tick style={pcTick, line width=0.4pt},
  },
  every axis plot/.append style={line width=0.9pt},
  /pgfplots/error bars/error bar style={line width=0.5pt},
}

% ---- 色板循环列表 ----
\pgfplotscreateplotcyclelist{pc academic}{
  pcBlue\\ pcOrange\\ pcGreen\\ pcMagenta\\ pcYellow\\ pcCyan\\ pcGray\\
}

% ---- 用法示例 ----
% \begin{tikzpicture}
%   \begin{axis}[cycle list name=pc academic, xlabel={年份}, ylabel={销量}, title={标题}]
%     \addplot coordinates {(1,2) (2,4) (3,3)};
%   \end{axis}
% \end{tikzpicture}

% ============================================================
% business 档（T2）差异：字号放大一档（label \normalsize / tick \small），
% cycle list 换 4E79A7/F28E2B/E15759/76B7B2/59A14F/EDC948/B07AA1/FF9DA7/9C755F/BAB0AC
% showcase 档（T3）：字号再放大（label \large / tick \normalsize），line width=1.4pt，
% cycle list 换 0077BB/EE7733/009988/EE3377/33BBEE/CC3311/BBBBBB
% ============================================================
，再 \input 本文件，正文选模式后作图。
% 依赖库已在此加载：shapes.geometric, arrows.meta, positioning
% ============================================================
\usetikzlibrary{shapes.geometric, arrows.meta, positioning}

% ---- 论文版（黑白无底色；字体由文档设：中文宋体 + 西文 Times New Roman）----
\newcommand{\pcPaperMode}{%
  \def\pcNodeLine{black!75}\def\pcNodeFill{white}\def\pcNodeText{black}%
  \def\pcKeyFill{white}\def\pcKeyLine{black}\def\pcKeyText{black}%
  \def\pcKeyWidth{1pt}% 关键节点靠粗边框区分，不用灰底
  \def\pcSubLine{black!40}\def\pcSubText{black!65}%
  \def\pcArrow{black!80}\def\pcLabel{black!70}%
  \def\pcNodeFont{\small}\def\pcFlowWidth{0.6pt}\def\pcStealth{2.2mm}}

% ---- PPT 版（showcase 色板：蓝主色、橙强调；字体建议中文黑体 + 西文 Arial）----
\newcommand{\pcPPTMode}{%
  \definecolor{pcPptBlue}{HTML}{0077BB}%
  \definecolor{pcPptOrange}{HTML}{EE7733}%
  \def\pcNodeLine{pcPptBlue!80}\def\pcNodeFill{pcPptBlue!10}\def\pcNodeText{black}%
  \def\pcKeyFill{pcPptBlue}\def\pcKeyLine{pcPptBlue}\def\pcKeyText{white}%
  \def\pcKeyWidth{\pcFlowWidth}%
  \def\pcSubLine{pcPptBlue!40}\def\pcSubText{black!65}%
  \def\pcArrow{black!70}\def\pcLabel{pcPptOrange!90!black}%
  \def\pcNodeFont{\normalsize}\def\pcFlowWidth{1.1pt}\def\pcStealth{3mm}}

\pcPaperMode % 默认论文版；PPT 场景在文档里改调 \pcPPTMode

% ---- 节点形状（形状语义与 references/diagram/flowchart.md 一致）----
\tikzset{
  % 起止节点：圆角胶囊
  pc start/.style={draw=\pcNodeLine, fill=\pcNodeFill, rounded corners=10pt,
                   minimum height=0.72cm, inner sep=7pt,
                   font=\pcNodeFont, text=\pcNodeText, align=center},
  % 处理：直角矩形
  pc process/.style={draw=\pcNodeLine, fill=\pcNodeFill, rounded corners=1.5pt,
                     minimum width=2.2cm, minimum height=0.75cm, inner sep=5pt,
                     font=\pcNodeFont, text=\pcNodeText, align=center},
  % 关键节点（关键路径）：论文版加粗边框（无底色），PPT 版主色实心白字
  pc key/.style={pc process, draw=\pcKeyLine, fill=\pcKeyFill, text=\pcKeyText,
                 line width=\pcKeyWidth, font=\pcNodeFont},
  % 判断：菱形
  pc decision/.style={shape=diamond, aspect=1.7, inner sep=1.5pt,
                      draw=\pcNodeLine, fill=\pcNodeFill,
                      font=\pcNodeFont, text=\pcNodeText, align=center},
  % 输入/输出：平行四边形
  pc io/.style={shape=trapezium, trapezium left angle=75, trapezium right angle=105,
                inner sep=4pt, minimum height=0.72cm,
                draw=\pcNodeLine, fill=\pcNodeFill,
                font=\pcNodeFont, text=\pcNodeText, align=center},
  % 数据库：圆柱
  pc db/.style={shape=cylinder, shape border rotate=90, cylinder uses custom fill,
                aspect=0.35, inner sep=2pt, minimum width=1.9cm, minimum height=1.0cm,
                draw=\pcNodeLine, fill=\pcNodeFill,
                font=\pcNodeFont, text=\pcNodeText, align=center},
  % 阶段/边界容器：虚线框
  pc sub/.style={draw=\pcSubLine, dashed, rounded corners=3pt, inner sep=8pt},
  pc sub label/.style={font=\footnotesize, text=\pcSubText},
}

% ---- 箭头与标注 ----
\tikzset{
  pc flow/.style={-{Stealth[length=\pcStealth]}, line width=\pcFlowWidth, draw=\pcArrow},
  pc flow back/.style={pc flow, dashed},                 % 回流/反馈：虚线
  pc label/.style={font=\footnotesize, text=\pcLabel, fill=white, inner sep=1.5pt},
}

% ---- 用法示例 ----
% \begin{tikzpicture}[node distance=8mm and 10mm]
%   \node[pc start]                (s) {开始};
%   \node[pc process, below=of s]  (a) {参数校验};
%   \node[pc decision, below=of a] (b) {库存充足？};
%   \draw[pc flow] (s) -- (a);
%   \draw[pc flow] (a) -- node[pc label, right=2pt]{通过} (b);
% \end{tikzpicture}
